\uuid{aiZB}
\titre{Théorème des valeurs intermédiaires, fonctions bornées}
\auteur{bodin}
\organisation{exo7}
\licence{CC BY-NC-SA 4.0}
\datecreate{2020-01-21}
\niveau{L1}
\module{Analyse}
\chapitre{Continuité, limite et étude de fonctions réelles}
\sousChapitre{Continuité : théorie}
\difficulte{}
\theme{Réels}
\exo7qcmid{arnaud.arnaud.026}

\contenu{
\begin{question}
\qclassification{123.01}
Quelles sont les assertions vraies ?
\begin{answers}
    \good{La méthode de dichotomie est basée sur le théorème des valeurs intermédiaires.}
    \bad{Tout polynôme de degré $\ge 3$ a au moins une racine réelle.}
    \bad{La fonction $f(x) = \frac{1}{x^3(x^2+1)}$ admet au moins une racine réelle dans $]-1,+1[$.}
    \good{Pour $f : \Rr^+ \longrightarrow \Rr$ continue admettant une limite finie en $+\infty$, $f$ est bornée.}
    \bad{Pour $f : \Rr^+ \longrightarrow \Rr$ continue admettant une limite finie qui vaut $f(0)$ en $+\infty$ alors $f$ est bornée et atteint ses bornes.}
\end{answers}
\end{question}
}
